\section{Implementation Validation}
\label{app:validation}

% No counterpart in Wu & Tu. It exists because two of these checks test the derivation rather
% than the code, and because a causal model that is silently non-causal produces entirely
% plausible numbers.

\TODO{the full checklist, with pass/fail and the numbers. \Cref{sec:exp-correctness} reports
only the two that test the equations; these are the rest.}

\begin{itemize}
  \item \textbf{Causality.} Changing the token at position $t$ leaves the logits at every
        position $<t$ bitwise unchanged (CPU, fixed seed --- on GPU, non-deterministic
        reductions can perturb the last bit of correct code).
  \item \textbf{Frozen prefix.} Prefix quantities are kept as attached leaves with
        \texttt{requires\_grad=True}; after backpropagating the loss at step $t$, their gradient
        is exactly zero for $j \ge t$ and non-zero for $j < t$. Testing this on a detached
        tensor would pass unconditionally and is therefore no test at all.
  \item \textbf{Tying.} Input and output embeddings are the same tensor object, not merely
        equal.
  \item \textbf{Normalisation.} All posteriors sum to one along their variable axis at every
        iteration.
  \item \textbf{Free energy.} The mean-field free energy is non-increasing across iterations on
        a toy graph. \TODO{figure} \NOTE{this is the check that catches a typo in the update
        equations; shapes, normalisation and causality are all blind to one.}
  \item \textbf{Exact readout vs.\ brute force.} On one slot the graph is a tree, so
        \Cref{eq:exact-readout} must agree with explicit enumeration over a tiny vocabulary.
        \TODO{agreement to 1e-N} This validates the factorisation itself.
  \item \textbf{Single-batch overfit.} Loss on one batch falls to near zero.
\end{itemize}
